\documentclass{article}
\begin{document}
\section{Esperimenti sul pendolo}
\subsection{Introduzione}
Questo documento è una relazione all'esperimento svolto a scuola utilizzando il pendolo semplice.
\subsection{Obiettivi: CONCLUDERE}
Gli obiettivi degli esperimenti svolti con il pendolo sono principalmente due. \\
Il primo è quello di calcolare l'accelerazione gravitazionale, $g$. Abbiamo infatti dimostrato in classe che con piccole oscillazioni del pendolo l'accelerazione gravitazionale è facilmente approssimabile. Nell'esperimento si è cercato di tenere angoli iniziali piccoli, per fare in modo che l'approssimazione abbia un basso errore. \\
Il secondo obiettivo è quello di verificare l'indipendenza del periodo d'oscillazione dal valore della massa che viene utilizzata. Sempre dalla formula dimostrata in classe si può notare come il periodo non doverebbe dipendere dalla massa (sempre solo con piccole oscillazioni).
\subsection{Teoria}
La formula dimostrata in classe è la seguente.
$$T = 2\pi\sqrt\frac{L}{g}$$
dove $T [\mathrm{s}]$ è il periodo d'oscillazione, dunque il tempo trascorso affinché la massa ritorni alla posizione iniziale, $\pi$ è il rapporto tra circonferenza e diametro e vale sempre $\pi = 3.1415926535897...$. Per i calcoli tuttavia si utlizzerà $\pi = 3.142$. $L [\mathrm{m}]$ è la lunghezza del filo utilizzato e $g [\mathrm{m s^{-2}}]$ è l'accelerazione gravitazionale, che calcoleremo utilizzando l'equazione di gravitazione universale, e l'altezza sul livello del mare del laboratorio in cui si è svolto l'esperimento. \\
Bisogna notare che questa formula si dimostra in ipotesi che l'angolo iniziale sia sufficentemente piccolo ($\theta_0 \approx 5^{\circ}$) e trascurando l'attrito dell'aria
Dalla formula citata si può ricavare $g$ con alcuni passaggi algebrici
$$T^2 = 4\pi^2\frac{L}{g}$$
si può elevare al quadrato perché entrambi i lati contengono valori positivi; il periodo per convenzione è sempre positivo e una radice quadrata non può essere negativa
$$g T^2 = 4\pi^2 L$$
$$g = 4\pi^2\frac{L}{T^2}$$
Questa è formula che si utilizzerà per calcolare l'accelerazione gravitazionale, dato un periodo e una lunghezza misurati. \\
Si calcola però anche la cosiddetta $g$ teorica, per poi computare l'errore fra quella teorica e quella sperimentale, cioè quella calcolata con i parametri del pendolo. \\
L'accelerazione gravitazionale teorica si calcola utilizzando l'equazione di gravitazione universale di Newton
$$F = G\frac{M_T m_p}{r^2}$$
dove $F [N]$ è la forza di attrazione gravitazionale tra il pendolo e la terra (ovviamente le altre attrazioni sono trascurabili), $G$ è la costante di gravitazione universale, che vale circa $G \approx 6.67 \cdot 10^{-11}\mathrm{N m^2 kg ^ {-2}}$, $M_T [\mathrm{kg}]$ è la massa della terra, $M_p [\mathrm{kg}]$ è la massa del pendolo e $r [\mathrm{m}]$ è la distanza tra il pendolo e il centro della terra, data da
$$r = r_t + h$$
dove $r_t \approx 6.373 \cdot 10^6[\mathrm{m}]$ è il raggio della terra e $h [\mathrm{m}]$ è l'altezza del pendolo sul livello del mare.
\subsection{Strumenti utilizzati}
Gli strumenti utilizzati sono:
\begin{itemize}
	\item sostegno per reggere il pendolo
	\item corda lunga (circa di $50 \mathrm{cm}$)
	\item corda corta (circa di $30 \mathrm{cm}$)
	\item sferetta metallica di colore dorato di massa $m = (72.3 \pm 0.1) \mathrm{g}$
	\item sferetta metallica di colore argenteo di massa $m = (24.4g \pm 0.1) \mathrm{g}$
	\item bilancia digitale di sensibilità $\pm 0.1 \mathrm{g}$
	\item metro a nastro
	\item goniometro
	\item cronometro digitale
	\item orologio digitale con sensore di altitudine
\end{itemize}
\subsection{descrizione}
L'esperienza è costituita da cinque esperimenti analoghi, che si differenziano solamente per il tipo di ogni strumento utilizzato. \\
Inizialmente si hanno misurato le masse delle due sferette.
In generale ogni esperimento si è svolto prendendo la corda e legandola per un estremità alla massa con un nodo e per l'altra al capo del sostegno. Si misura quindi la lunghezza tra il l'origine della corda nel sostegno e il centro della massa. Dopodichè si prende la massa e si sposta tenendo la corda tesa tale da formare un angolo iniziale, che si misura con il goniometro. A questo punto di lascia la massa facendo partire nello stesso istante il cronometro. Si lascia che la massa compia dieci oscillazioni e si ferma il cronometro. Nel mio caso si ha misurato il tempo trascorso tra dieci oscillazioni con due cronometri.\\
Detto ciò, seguono i dati ricavati da ogni esperimento.

\subsubsection{Esperimento 1}
sferetta dorata di massa $m_d = (72.3 \pm 0.1) \mathrm{g}$\\
lunghezza $L = (46.8 \pm 0.2) \mathrm{cm}$\\
angolo iniziale $\theta_0 = (8 \pm  0.2)^{\circ}$\\
tempo misurato dal primo cronometro $t_1 = 13.16 \mathrm{s}$\\
tempo misurato dal secondo cronometro $t_2 = 12.95 \mathrm{s}$

\subsubsection{Esperimento 2}
sferetta argentea di massa $m_a = (24.4 \pm 0.1) \mathrm{g}$\\
lunghezza $L = (40.5 \pm 0.2) \mathrm{cm}$\\
angolo iniziale $\theta_0 = (11 \pm 0.1)^{\circ}$\\
tempo misurato dal primo cronometro $t_1 = 12.32 \mathrm{s}$\\
tempo misurato dal secondo cronometro $t_2 = 12.15 \mathrm{s}$

\subsubsection{Esperimento 3}
sferetta argentea di massa $m_a = (24.4 \pm 0.1) \mathrm{g}$\\
lunghezza $L = (41.5 \pm 0.2) \mathrm{cm}$\\
angolo iniziale $\theta_0 = (7 \pm 0.1)^{\circ}$\\
tempo misurato dal primo cronometro $t_1 = 12.23 \mathrm{s}$\\
tempo misurato dal secondo cronometro $t_2 = 12.07 \mathrm{s}$

\subsubsection{Esperimento 4}
sferetta dorata di massa $m_d = (72.3 \pm 0.1) \mathrm{g}$\\
lunghezza $L = (41.5 \pm 0.2) \mathrm{cm}$\\
angolo iniziale $\theta_0 = (7 \pm 0.1)^{\circ}$\\
tempo misurato dal primo cronometro $t_1 = 12.32 \mathrm{s}$\\
tempo misurato dal secondo cronometro $t_2 = 12.13 \mathrm{s}$

\subsubsection{Esperimento 5}
sferetta argentea di massa $m_a = (24.4 \pm 0.1) \mathrm{g}$\\
lunghezza $L = (27.8 \pm 0.2) \mathrm{cm}$\\
angolo iniziale $\theta_0 = (5 \pm 0.1)^{\circ}$\\
tempo misurato dal primo cronometro $t_1 = 9.66 \mathrm{s}$\\
tempo misurato dal secondo cronometro $t_2 = 9.76 \mathrm{s}$

\end{document}
